\begin{abstract}
\noindent
Understanding how carbon flows through distinct soil channels is critical for predicting ecosystem resilience. 
Here, we present a compartmental modeling framework to analyze the coupling of litter-derived (Synapse) 
and root-derived (Vertex) carbon pathways in terrestrial food webs. Using a multi-channel differential model, 
we simulate how microbial (Nimbus) search behavior and predation by Apex consumers regulate trophic cascades 
under climate fluctuations. Our model shows that asynchronous resource inputs across the two channels stabilize 
macro-consumer biomass by preventing concurrent resource crashes. However, high-frequency rainfall extremes 
decouple the channels, causing transient synchronization that increases extinction risk for top predators by 34\%. 
These results demonstrate that resource pathway diversity act as a crucial buffer for food web stability. 
We discuss how these findings can guide soil conservation strategies in vulnerable grassland ecosystems.
\end{abstract}

\vspace{0.5cm}
\noindent \textbf{Keywords:} Carbon cycling, Compartmental model, Ecology Letters, Food web stability, Multi-channel pathways, Soil food web.

\begin{metadata-box}
\begin{spacing}{1.1}
\sffamily\small
\begin{tabularx}{\textwidth}{lX}
\textbf{Article Type:} & Letter \\
\textbf{Word Count (Abstract):} & 142 words (Max 150 words) \\
\textbf{Word Count (Main Text):} & 4,820 words (Max 5,000 words) \\
\textbf{Number of References:} & 28 \\
\textbf{Number of Figures / Tables:} & 2 Figures / 1 Table \\
\textbf{Corresponding Author:} & Dr. Arthur Vance (\href{mailto:a.vance@vertex-univ.edu}{\texttt{a.vance@vertex-univ.edu}}) \\
\textbf{Submission Target:} & \textit{Ecology Letters} \\
\end{tabularx}
\end{spacing}
\end{metadata-box}
